\chapter{Arquitetura da rota de interpolação nativa}
\label{cap:arquitetura}

\section{A rota original e sua limitação estrutural}
\label{sec:rota_original}

O \selfgrib{} original gera condição inicial e fronteira lateral em três
estágios, ilustrados na metade superior da \cref{fig:comparacao_rotas}:
\texttt{extract\_fields} interpola verticalmente os campos nativos do
\mpasa{} (55 níveis de modelo) para um conjunto de 61 níveis de pressão
fixa, ainda sobre a malha nativa \emph{global}; \texttt{convert\_mpas}
(ferramenta da NCAR, reaproveitada sem modificação) remapeia
horizontalmente esses campos da malha nativa para uma grade
latitude-longitude regular; e \texttt{pack\_intermediate} escreve o
resultado no formato binário intermediário do WPS. Esse arquivo
intermediário é então consumido pelo \texttt{init\_atmosphere\_model}
original (inalterado), que interpola de volta -- bilinear ou 16 pontos,
sobre índices de grade estruturada -- para a malha nativa \emph{regional}
de destino, e resolve internamente a grade vertical, a interpolação
vertical e o balanço hidrostático.

Como discutido na \cref{cap:introducao}, essa rota tem uma limitação
estrutural: o \texttt{init\_atmosphere\_model} não aceita, para o caso de
dado real, nenhuma forma de entrada que não seja uma grade estruturada
-- confirmado pela leitura da rotina \texttt{mpas\_init\_atm\_hinterp.F}
do código-fonte de referência. Isso significa que o ganho de precisão de
uma interpolação nativa malha-a-malha (\cref{sec:baricentrica}) só é
alcançável substituindo \emph{inteiramente} o caminho
\texttt{convert\_mpas}~$\to$~formato WPS~$\to$~\texttt{init\_atmosphere\_model},
não apenas uma parte dele.

\section{A rota nativa: visão geral}
\label{sec:rota_nativa}

A rota descrita neste relatório (metade inferior da
\cref{fig:comparacao_rotas}) substitui esse caminho por sete fases,
implementadas como programas Fortran independentes que se comunicam por
arquivos NetCDF intermediários -- mantendo a mesma filosofia de
composição em estágios discretos e idempotentes já usada na rota
original, mas eliminando toda e qualquer grade regular do caminho de
dados.

\begin{figure}[H]
  \centering
  \begin{tikzpicture}[
      node distance=6mm and 10mm,
      box/.style={draw, rounded corners, minimum height=0.9cm, minimum width=2.6cm, align=center, font=\footnotesize, fill=blue!6},
      grid/.style={draw, minimum height=0.9cm, minimum width=2.6cm, align=center, font=\footnotesize, fill=red!8, dashed},
      lbl/.style={font=\scriptsize\itshape, text width=2.6cm, align=center},
      arr/.style={-{Stealth[length=2mm]}, thick},
    ]

    % ===================== ROTA ORIGINAL =====================
    \node[font=\bfseries\small] (titulo1) at (0,4.6) {Rota original (via WPS)};

    \node[box] (g1) at (-6.5,3.6) {malha nativa\\GLOBAL};
    \node[box] (e1) at (-3.7,3.6) {\texttt{extract\_}\\\texttt{fields}};
    \node[grid] (c1) at (-0.9,3.6) {grade\\lat-lon};
    \node[box] (p1) at (1.9,3.6)  {formato\\WPS};
    \node[box] (i1) at (4.9,3.6)  {\texttt{init\_atmosphere\_}\\\texttt{model} (original)};
    \node[box] (r1) at (7.9,3.6)  {malha nativa\\REGIONAL};

    \draw[arr] (g1) -- (e1);
    \draw[arr] (e1) -- (c1) node[midway, above, font=\tiny] {\texttt{convert\_mpas}};
    \draw[arr] (c1) -- (p1) node[midway, above, font=\tiny] {\texttt{pack\_intermediate}};
    \draw[arr] (p1) -- (i1);
    \draw[arr] (i1) -- (r1) node[midway, above, font=\tiny] {bilinear/16-pt};

    \node[lbl, red!50!black] at (-0.9,2.75) {\textbf{grade intermediária}\\(fonte do erro de\\ida-e-volta)};

    % ===================== ROTA NATIVA =====================
    \node[font=\bfseries\small] (titulo2) at (0,1.4) {Rota nativa (\emph{voronoi-to-voronoi}, este trabalho)};

    \node[box] (g2) at (-6.5,0.4) {malha nativa\\GLOBAL};
    \node[box, fill=green!10] (f1) at (-3.7,0.4) {Fase 1\\\texttt{hinterp\_native}};
    \node[box, fill=green!10] (f2) at (-0.9,0.4) {Fase 2\\grade vertical};
    \node[box, fill=green!10] (f34) at (1.9,0.4) {Fases 3+4\\interp. vert. +\\hidrostático};
    \node[box, fill=green!10] (f6) at (4.9,0.4) {Fase 6\\superfície/solo};
    \node[box] (r2) at (7.9,0.4) {\texttt{init.nc} /\\\texttt{lbc.*.nc}};

    \draw[arr] (g2) -- (f1);
    \draw[arr] (f1) -- (f2);
    \draw[arr] (f2) -- (f34);
    \draw[arr] (f34) -- (f6);
    \draw[arr] (f6) -- (r2);

    \node[box, fill=green!10] (f7) at (7.9,-0.9) {Fase 7\\\texttt{mpas\_atmosphere}};
    \draw[arr] (r2) -- (f7);

    \node[lbl, green!35!black] at (-3.7,-0.5) {interpolação\\baricêntrica\\direta malha--malha};

  \end{tikzpicture}
  \caption{Comparação entre a rota original (via formato binário WPS e
    grade lat-lon intermediária, acima) e a rota nativa
    \emph{voronoi-to-voronoi} apresentada neste trabalho (abaixo). A rota
    nativa elimina toda etapa de reamostragem em grade regular: da malha
    de Voronoi global de origem à malha de Voronoi regional de destino,
    o dado nunca deixa de estar sobre uma malha não estruturada.}
  \label{fig:comparacao_rotas}
\end{figure}

As sete fases, cujo detalhamento técnico ocupa os
Capítulos~\ref{cap:fase1}--\ref{cap:fase7}, são:

\begin{description}[leftmargin=2.2cm, style=nextline]
  \item[Fase 1 -- Interpolação horizontal nativa] Remapeia os campos já
    extraídos e verticalmente interpolados para níveis de pressão fixa
    (ainda na malha global) diretamente para os centros de célula da
    malha regional de destino, via interpolação baricêntrica
    (\cref{sec:baricentrica}), sem grade intermediária.
  \item[Fase 2 -- Grade vertical nativa] Calcula a geometria vertical
    \textit{terrain-following} híbrida da malha de destino
    ($z_{grid}$, $zz$, coeficientes de interpolação vertical, termos de
    métrica para advecção) a partir apenas do terreno e da conectividade
    dessa malha -- não depende do dado meteorológico.
  \item[Fase 3 -- Interpolação vertical] Interpola os campos já
    horizontalmente remapeados (ainda em níveis de pressão fixa) para os
    níveis de modelo da grade vertical nativa calculada na Fase~2.
  \item[Fase 4 -- Balanço hidrostático] Resolve o campo de pressão,
    densidade seca e temperatura potencial (úmida) de forma consistente
    com a métrica vertical, além de água precipitável e vento vertical
    diagnóstico.
  \item[Fase 6 -- Campos de superfície e solo] Calcula os campos que não
    vêm diretamente da coluna atmosférica: temperatura profunda do solo,
    correção de temperatura de pele por diferença de elevação, fração de
    vegetação/albedo de fundo por interpolação climatológica mensal,
    flag de neve/gelo marinho.
  \item[Fase 7 -- Escrita e execução] Escreve os arquivos
    \texttt{init.nc}/\texttt{lbc.*.nc} completos (mesclando os campos
    calculados nas fases anteriores com os campos de cópia direta da
    malha estática) e alimenta o núcleo \texttt{mpas\_atmosphere} real.
\end{description}

\begin{center}
\begin{minipage}{0.9\textwidth}
\small\itshape
Nota de numeração: a Fase~5 do plano de implementação original (projeção
de vento células$\to$arestas) foi absorvida dentro das Fases~3/4, já que
a reconstrução da componente normal do vento e sua multiplicação pela
densidade ($ru$) fazem parte, respectivamente, da interpolação vertical
e do balanço hidrostático -- não há um programa Fase~5 isolado.
\end{minipage}
\end{center}

\section{Decisão de escopo: por que não um subconjunto}
\label{sec:decisao_escopo}

Uma primeira hipótese de trabalho, descartada ainda na fase de
investigação, era um ``escopo mínimo'': gerar o dado meteorológico em
modo \emph{scatter} (lista de pontos exatos da malha regional, em vez de
grade), mas ainda escrever no formato binário WPS e deixar o
\texttt{init\_atmosphere\_model} original consumi-lo. A leitura do
código-fonte de \texttt{mpas\_init\_atm\_hinterp.F} mostrou que essa
hipótese está estruturalmente quebrada: o formato WPS \emph{é}, por
definição, uma grade regular (o próprio cabeçalho de cada registro do
formato carrega \texttt{nx}, \texttt{ny}, espaçamento e origem de uma
grade); não existe uma forma de ``passar direto'' pontos dispersos por
esse caminho. Mesmo alimentando o \texttt{convert\_mpas} em modo
\emph{scatter} para gerar um arquivo tecnicamente válido, o
\texttt{init\_atmosphere\_model} ainda executaria um segundo salto de
reinterpolação bilinear/16 pontos ao final -- o ganho de precisão
validado no protótipo inicial (recuperação exata nos centros de célula,
\cref{sec:validacao_fase1}) só é alcançável pulando esse caminho por
completo, não reduzindo seu escopo.

Por essa razão, a decisão de projeto foi pelo escopo completo: reescrever
cada etapa fisicamente necessária -- grade vertical, interpolação
vertical, balanço hidrostático, campos de superfície -- como módulos
Fortran novos, extraídos literalmente do código-fonte de referência do
MPAS-Model (\cref{cap:implementacao}), em vez de reaproveitar o
\texttt{init\_atmosphere\_model} para qualquer parte dessas etapas.

\section{Generalização: leitura do \textit{namelist} real em tempo de execução}
\label{sec:generalizacao}

Uma decisão de engenharia transversal a todas as fases, adotada após
revisão explícita do escopo do trabalho, foi banir do código-fonte
qualquer parâmetro de configuração (\texttt{config\_*}) fixo como
constante de compilação. Nas primeiras versões dos programas, valores
como \texttt{config\_ztop}, \texttt{config\_dzmin},
\texttt{config\_nsm}, e -- criticamente -- \texttt{config\_start\_time}
(a própria data/hora da rodada) estavam hard-coded, o que restringia a
validação a um único caso de teste e, mais grave, impedia diretamente a
geração da fronteira lateral (\texttt{lbc.*.nc}), já que cada arquivo de
fronteira corresponde a um tempo diferente dentro do mesmo experimento.
Um módulo dedicado (\texttt{namelist\_config.F90},
\cref{sec:namelist_config}) foi introduzido para ler o
\texttt{namelist.init\_atmosphere} \emph{real} do experimento, via
\texttt{NAMELIST} nativo do Fortran, em tempo de execução -- o mesmo
arquivo de configuração já usado pela rota original, sem introduzir
nenhum formato novo. Essa mudança, além de generalizar a ferramenta para
qualquer malha/experimento, revelou um bug real de suposição incorreta
sobre o valor padrão de \texttt{config\_frac\_seaice}
(\cref{cap:discussao}), que só se manifestaria em malhas com gelo
marinho -- não presente no caso de validação usado neste relatório, mas
corrigido preventivamente.
